\section{A complete stationary response experiment}
\label{s4:response-construction}
The moment identity in Equation~\eqref{r3:lpiv-ratio} becomes a usable
operation only after specifying how observations repeat. Consider a
three-dimensional vector containing a policy rate, deposits and lending,
all in declared units. Let
\[
 X_t=FX_{t-1}+B\varepsilon_t,\qquad
 Z_t=a\varepsilon_{1t}+v_t,\qquad
 v_t=\rho_vv_{t-1}+\sqrt{1-\rho_v^2}\zeta_t.
\]
The innovations are mutually independent standard Gaussian vectors/scalars;
the joint process starts in its stationary law. The eigenvalues of $F$ and
$\rho_v$ have modulus below one, with a fixed stability margin. The matrix
$B$ is nonsingular, $B_{11}\ne0$, and initially $aB_{11}$ is bounded away
from zero. Gaussianity is a statistical restriction of this reference
experiment, while the instrument restrictions supply causal identification.
The Gaussian observational VAR itself does not identify $B$.

Fix $H$ before sampling, and use origins $t=1,\ldots,n=T-H$ for every
coordinate. Index $j=(i,h)$ for $i\in\{2,3\}$ and $0\le h\le H$.
Set $U_{tj}=Z_tX_{i,t+h}$, $D_t=Z_tX_{1t}$,
$q=\E D_t=aB_{11}$, and $\theta_j=\E U_{tj}/q$.
The estimate and its exact linearization are
\begin{equation}
 \widehat\theta_j=\bar U_j/\bar D,\qquad
 \sqrt n(\widehat\theta_j-\theta_j)
 =\frac{1}{\bar D}\frac1{\sqrt n}\sum_{t=1}^n
       Z_t(X_{i,t+h}-\theta_jX_{1t}).
 \label{s4:ratio-linearization}
\end{equation}
This identity avoids a hidden regression remainder. A one-percentage-point
impact normalization is meaningful only if the policy rate is expressed
in percentage points. A change to basis points rescales the response.

The score $M_{tj}=U_{tj}-\theta_jD_t$ has mean zero. The long-run matrix
$\Omega=\sum_{k\in\mathbb Z}\Cov(M_t,M_{t-k})$ contains both serial
and cross-horizon dependence. In particular, independence of $v$ and $X$
does not make their product serially independent. Its contribution to
the lag covariance is
$\Cov(v_t,v_{t-k})\Cov(X_{i,t+h}-\theta_jX_{1t},
X_{i',t-k+h'}-\theta_{j'}X_{1,t-k})$.
Persistent measurement noise in an otherwise valid instrument can therefore
matter greatly for standard errors. This is the score-level distinction
emphasized by \citet[Section 2.3]{stock2018}.

\subsection{Why a central limit theorem applies here}
For this particular experiment the limiting argument can be reconstructed
directly. The jointly Gaussian vector process $(X_t,Z_t)$ has geometrically
decaying autocovariances and a bounded spectral density. Gaussian moment
pairing shows that the covariance of any two centered quadratic products
is a sum of products of two covariances. Hence the score covariances are
absolutely summable; a fixed horizon shift changes the constants, not
summability.

For any fixed vector $c$, collect the finitely many Gaussian observations
needed for $\sum_tc'M_t$ into $G_n$. That sum is a centered quadratic
form $G_n'A_nG_n-\operatorname{tr}(A_n\Sigma_n)$. Each observation is
paired with only a fixed number of lags, so $\|A_n\|$ is bounded.
Absolute covariance summability bounds $\|\Sigma_n\|$ uniformly. The
eigenvalues $\lambda_{nr}$ of
$\Sigma_n^{1/2}A_n\Sigma_n^{1/2}$ are consequently bounded, and the form
has the law $\sum_r\lambda_{nr}(\chi_{1r}^2-1)$ with independent terms.
Its variance divided by $n$ converges to $c'\Omega c$. If this limit is
positive, the fourth-moment Lyapunov ratio is bounded by
$C\sum_r\lambda_{nr}^4/(\sum_r\lambda_{nr}^2)^2=O(n^{-1})$.
If the limit is zero, the normalized form converges to zero in $L_2$.
Applying this argument to every $c$ proves the joint Gaussian limit.
With $\bar D\to q$, Equation~\eqref{s4:ratio-linearization} gives
$\sqrt n(\widehat\theta-\theta)\Rightarrow N(0,\Omega/q^2)$.
This is a project proof for a fixed-dimensional Gaussian linear law;
stationarity alone does not prove its conclusion.

\subsection{A covariance estimator with an explicit bandwidth}
Let $\widehat M_t=U_t-\widehat\theta D_t$ and center each score column.
With $L=\lfloor n^{1/4}\rfloor$, compute
\begin{align}
 \widehat\Gamma_k&=n^{-1}\sum_{t=k+1}^n
    (\widehat M_t-\bar M)(\widehat M_{t-k}-\bar M)',\nonumber\\
 \widehat\Omega_L&=\widehat\Gamma_0+
   \sum_{k=1}^L\left(1-\frac{k}{L+1}\right)
       (\widehat\Gamma_k+\widehat\Gamma_k'),\qquad
 \widehat V=\frac{\widehat\Omega_L}{n\bar D^2}.
 \label{s4:bartlett}
\end{align}
The common denominator $n$, including at nonzero lags, matters. Extending
the centered sequence by zeros, this estimator is the sum of outer
products of all length-$(L+1)$ moving sums, divided by $n(L+1)$. It is
positive semidefinite. The original construction and a general consistency
result are in \citet[equation (5), Theorems 1--2]{newey1986}; the latter's
general mixing theorem requires $L=o(n^{1/4})$. Our stated $n^{1/4}$
choice instead uses the stronger Gaussian structure just specified.

Here are the additional steps. Gaussian moment pairing bounds the variance
of every empirical lag covariance by $C/n$, uniformly in the included
lag: in each covariance of two fourth-degree terms, a surviving pairing
joins the two date groups, whose geometrically decaying covariance is
summable over their separation. There are finitely many pairings, and
shifts by $k$ do not change that sum. Summing the resulting standard
deviation bounds gives $O_p(L/\sqrt n)$. The Bartlett taper bias is
bounded by $(L+1)^{-1}\sum_k |k|\|\Gamma_k\|$ and the omitted tail
is geometric. Replacing the true parameter by a root-$n$ estimate changes
the sum by $O_p(L/\sqrt n)$, using Cauchy--Schwarz on each empirical
cross product. Centering and edge corrections vanish as well. Thus
\[
 \|\widehat\Omega_L-\Omega\|
 =O_p(L/\sqrt n)+O(L^{-1})+O(r^L),\qquad r<1,
\]
is a conservative bound in this fixed model. It proves consistency for
the stated bandwidth; it does not promise small bias at 800 monthly origins
or uniformly as $r$ approaches one.

For $d=2(H+1)$ response coordinates and positive marginal variances,
$\widehat\theta_j\pm z_{1-\alpha/(2d)}\sqrt{\widehat V_{jj}}$
is an asymptotic simultaneous Bonferroni band. A less conservative
alternative computes the $(1-\alpha)$ quantile of
$\max_j|G_j|$ for $G\sim N(0,\widehat R)$, where $\widehat R$ is
the correlation matrix of $\widehat V$. Numerical Gaussian draws
approximate that quantile; their Monte Carlo error is additional to the
sampling approximation. The experiments use 4,096 such draws and report
that limitation. A null-centered significance band tests a specified
zero-response hypothesis and should not be relabeled as a confidence
band for every response; see \citet[equations (5)--(8)]{inoue2023}.

\subsection{Weak instruments require a different result object}
Do not divide by a weak normalization to construct the inferential pivot.
For a proposed response $b$, test the zero mean of $U_{tj}-bD_t$.
Write $S_{UU},S_{UD},S_{DD}$ for the corresponding entries of the
Bartlett matrix of the centered joint products $(U_t,D_t)$. Accept $b$ if
\begin{equation}
 n(\bar U_j-b\bar D)^2
 \le c^2(S_{UU}-2bS_{UD}+b^2S_{DD}).
 \label{s4:weak-inversion}
\end{equation}
This is one quadratic inequality $A b^2+B b+C\le0$, with
$A=n\bar D^2-c^2S_{DD}$,
$B=-2n\bar U_j\bar D+2c^2S_{UD}$, and
$C=n\bar U_j^2-c^2S_{UU}$.
Its solution can be an interval, two unbounded rays, the whole line,
an empty set, or a degenerate singleton/half-line. These cases must survive
the interface. At the true $b$, the score CLT does not require dividing
by $q$, and remains applicable along bounded weak-instrument sequences
with a nondegenerate residual-score variance and fixed stability margin.
Bonferroni inversion across coordinates yields simultaneous coverage of
the true response vector. At $q=0$ the normalized causal response is
not identified by these moments; a statistical set cannot repair that.
No first-stage threshold is used to switch silently between guarantees.

\subsection{A fully specified propagation challenger}
Fit the VAR(1) by OLS, giving $\widehat F$, residuals $\widehat u_t$ and
$\widehat\kappa=n^{-1}\sum_tZ_t\widehat u_t$. Its response is
$e_i'\widehat F^h\widehat\kappa/\widehat\kappa_1$.
With $Q=\E X_{t-1}X_{t-1}'$, the influence matrix for $F$ is
$u_tX_{t-1}'Q^{-1}$. The influence vector for $\kappa$ is
$Z_tu_t-\kappa$ because the stipulated exclusions give
$\E[Z_tX_{t-1}]=0$; without that condition an additional derivative term
is required. For perturbations $(dF,d\kappa)$,
\[
 d\theta_{ih}=\frac{e_i'\sum_{r=0}^{h-1}F^r(dF)F^{h-1-r}\kappa}
                       {\kappa_1}
 +\frac{e_i'F^hd\kappa-\theta_{ih}d\kappa_1}{\kappa_1}.
\]
Apply this derivative to the joint score and use its joint HAC matrix.
This challenger propagates uncertainty in the fitted dynamics and impact
direction. Its dimensional savings rely on correct finite-order dynamics;
a misspecified VAR cannot borrow the direct moment estimator's robustness.

\section{Learned controls under explicit temporal separation}
\label{s4:dependent-learning}
Many bank analyses require hundreds of release-vintage controls. A
partially linear IV model supplies a precise learning target:
\[
 Y_t=\theta D_t+g(C_t)+e_t,\quad
 \E[e_t\mid C_t,Z_t]=0,\quad
 (\ell,r,m)(c)=\E[(Y_t,D_t,Z_t)\mid C_t=c].
\]
The orthogonal score and its identifying slope are
\[
 \psi_t(\theta,\eta)=(Z_t-m(C_t))
  \{Y_t-\ell(C_t)-\theta[D_t-r(C_t)]\},\qquad
 J=\E[(Z_t-m)(D_t-r)]\ne0.
\]
Instruments, controls and the outcome horizon must already satisfy the
causal restrictions. Good nuisance prediction does not supply them.
For fitted errors $\Delta m,\Delta\ell,\Delta r$, conditional on an
independent training sample, direct expansion gives the exact mean remainder
\begin{equation}
 \E\{\psi(\theta,\widehat\eta)-\psi(\theta,\eta)\}
 =\E[\Delta m(\Delta\ell-\theta\Delta r)].
 \label{s4:dependent-remainder}
\end{equation}
The first-order terms have zero means by their respective conditional
moment restrictions. This connects the dynamic procedure to the
orthogonality argument in Chapter~\ref{r3:learning} and
\citet{dml2018}, with a new sampling argument still needed.

One sufficient construction is deliberately explicit. Suppose the underlying
observables and fitted functions are bounded, the process is absolutely
regular with $\beta(k)\le C\exp(-ck)$, and lag/horizon windows are fixed.
Use two contiguous blocks of comparable size, separated by
$g=\lceil3\log n/c\rceil$ \emph{after} expanding both blocks for every
feature lag and future outcome. Train on either block and evaluate on the
other. Absolute regularity couples the pair of blocks to independent
blocks with error at most $C\exp(-cg)=O(n^{-3})$. Both fits use this
same coupled pair; independence of the two fitted estimates is not asserted.
An unknown $c$ is a scientific tuning uncertainty. The formula is a
sufficient regime specification, not an estimate of $c$ from a few scores.

Let $\delta_n$ bound the three population $L_2$ nuisance errors. For a
bounded centered score difference its lag covariance is bounded both by
$C\delta_n^2$ and by $C\exp(-ck)$. Splitting the lag sum at
$k\asymp\log(1/\delta_n)$ bounds its long-run variance by
$C\delta_n^2\{1+\log(1/\delta_n)\}$. Conditional centering, Chebyshev's
inequality and Equation~\eqref{s4:dependent-remainder} then give the oracle
linearization when
\[
 \delta_n\to0,\qquad
 \|\Delta m\|_2(\|\Delta\ell\|_2+|\theta|\|\Delta r\|_2)
       =o_p(n^{-1/2}).
\]
The oracle CLT can also be constructed here. Split the bounded centered
score into big blocks of length $m=\lfloor n^{1/3}\rfloor$, separated
by gaps $b=\lceil C\log n\rceil$ with $cC>2$. Summable score
covariances bound the variance of the discarded gaps by a constant times
their number of observations, $O(n^{2/3}\log n)=o(n)$. Coupling the
remaining big blocks to independent blocks fails with probability at most
$O(n^{2/3}e^{-cb})=o(1)$. Each centered block is bounded by $O(m)$,
which is $o(\sqrt n)$; hence the independent-block Lindeberg condition
holds. Their variance divided by $n$ converges to the summable long-run
variance. Cram\'er--Wold and the negligible gaps give the vector CLT.

For covariance estimation, a covariance of two score products decomposes
into a fourth cumulant plus two products of pairwise covariances. Bounded
exponential absolute regularity bounds the cumulant by an exponential in
the largest separating gap of the four dates; summing its possible gaps
is finite. Pairwise terms are summable as well. This gives a uniform
$C/n$ variance bound for empirical lag covariances, and the same
$O_p(L/\sqrt n)$ conservative stochastic bound as above. Across the deleted central gap,
do not pretend the two observed endpoints are adjacent. Sum lag products
within each evaluation block using the total retained length as denominator.
The missing oracle boundary pairs contribute $O(L^2/n)$, which vanishes
for $L=\lfloor n^{1/4}\rfloor$.

There is an additional condition for fitted-score covariance. If the
empirical root-mean-square score error is $r_n$, Cauchy--Schwarz bounds
the difference of the two HAC matrices by $O_p(Lr_n+Lr_n^2)$.
Take $\delta_n$ to be a deterministic high-probability bound on the
population errors. Conditional evaluation and Markov's inequality applied
to the mean squared score error give $r_n=O_p(\delta_n)$; the coupling
failure probability vanishes. Require $L\delta_n=o(1)$, as well as the product condition.
Taking all nuisance errors $o_p(n^{-1/4})$ is sufficient for the stated
bandwidth. A maximum lag-covariance error $o_p(1)$ alone cannot be
summed over an increasing number of lags.

\subsection{When sparsity supplies the required rates}
For a concrete learner suppose each conditional mean is a linear combination
of $p\le n^a$ bounded, standardized features, with at most $s$ nonzero
coefficients, plus an explicitly controlled approximation error. Assume
a population restricted eigenvalue bounded below. Partition training data
into alternating blocks of length $b=C_0\log n$, with
$C_0>(3+2a)/c$. Coupling alternate blocks, followed by bounded-block
concentration and a union bound over scores and Gram entries, gives
deviations of order $\lambda_n=C_1\sqrt{\log n\log(pn)/n}$ for a
constant $C_1$ depending on the envelope and dependence constants.
This conservative penalty includes a logarithmic dependence cost.

The Lasso basic inequality on the event
$\|n^{-1}\sum C_te_t\|_\infty\le\lambda_n/2$ restricts its error
to a cone. Restricted curvature then gives prediction error
$O_p(\sqrt{s}\lambda_n)$ and coefficient $\ell_1$ error
$O_p(s\lambda_n)$. Concentration of the Gram matrix transfers the
empirical prediction bound to population $L_2$ when $s\lambda_n=o(1)$.
Clipping predictions to the known outcome envelope cannot increase that
population error. Consequently
$s\log n\log(pn)=o(\sqrt n)$, with approximation errors obeying the
same product and bandwidth restrictions, supplies a complete sufficient
learning regime. This is substantially stronger than merely $p<n$.

\citet[Assumptions 3.1--3.4 and Theorem 3.1]{babii2020} develops a more
flexible sparse-group route under $\tau$-mixing and polynomial tails,
with a penalty depending explicitly on moments and decay rates. Its
prediction inequalities do not automatically supply the population nuisance
rates or the bounded coupling argument just used. A factor alternative
$C_t=\Lambda f_t+e_t^C$ economizes on a dense common component, but
dropping $e_t^C$ is valid adjustment only if the omitted component has no
remaining confounding role, or its bias is bounded. Sparsity and factor
reduction are therefore different scientific restrictions. Forecast RMSE
does not decide between them for causal inference.

The current proposal of \citet[Assumptions 2.1--2.5]{ciganovic2026}
uses reverse cross-fitting for time series and explicitly assumes
conditional plug-in stability. That assumption cannot be deduced from
disjoint row indices or reversibility. For example, a stationary Gaussian
process can be reversible and still have nonsummable long-memory
covariances. Its Supplement S1's increasing-bandwidth argument also needs
control of the accumulated fitted-score discrepancy. We retain the source
as a relevant 2026 alternative and use the stronger explicit separation
and rate construction above for the initial justified route. Removing
gaps requires a learner-specific stability proof.

\section{Common bank shocks determine the effective sample size}
\label{s4:date-panel}
For a fixed set of banks with target weights $w_i$, the relevant score at
date $t$ is $\bar\psi_t=\sum_iw_i\psi_{it}$, not an array of $NT$
independent rows. Apply the temporal construction to this aggregate when
its own conditions hold. For the transparent example
$\psi_{it}=C_t+e_{it}$, with independent idiosyncratic white noise,
\[
 \operatorname{avar}(\widehat\theta)
 =J^{-2}T^{-1}\left\{\Omega_C+\sigma_e^2\sum_iw_i^2\right\}.
\]
Equal weights reduce the idiosyncratic term by $N$ and leave the common
component unchanged. More banks do not create more monetary-policy shocks.
Sampling a random population of banks adds a separate uncertainty source;
time aggregation alone does not cover it. Staggered-treatment comparisons,
dynamic-panel moment restrictions and multiway-cluster methods from
Chapter~\ref{r3:panels} answer distinct questions and retain their own
assignment and asymptotic conditions.

\section{Persistence, cointegration and the level target}
\label{s4:integrated-responses}
An integrated balance series is economically plausible and falls outside
the stable-level experiment. For a VAR(1) with semisimple unit roots and
stable remaining roots,
\[
 \Delta X_t=(F-I)X_{t-1}+B\varepsilon_t
             =\alpha\beta'X_{t-1}+B\varepsilon_t.
\]
The rank of $F-I$ and the stationary error-correction combinations
$\beta'X_t$ describe long-run comovement. They do not identify which
innovation is a monetary, liquidity or credit-supply shock. The latter
still requires an instrument or another economic restriction.

There is an identified stationary transformation for the fixed-horizon
external-IV target. Suppose $\Delta X$ has a stationary Gaussian linear
representation with geometrically decaying coefficients and the original
lead/lag exclusions hold. Write $\theta_{-1}=0$. Then
\[
 \frac{\E[Z_t\Delta X_{i,t+h}]}{\E[Z_t\Delta X_{1t}]}
       =\theta_{ih}-\theta_{i,h-1},\qquad
 \theta_{ih}=\sum_{r=0}^h
       \frac{\E[Z_t\Delta X_{i,t+r}]}{\E[Z_t\Delta X_{1t}]}.
\]
The denominator still equals $aB_{11}$. The numerator sums to
$\E[Z_t(X_{i,t+h}-X_{i,t-1})]$: this is the response of the
\emph{level} to a one-time shock, recovered from stationary changes,
not the effect of permanently changing a rate rule. Estimate the stacked
change responses and premultiply by the block-diagonal cumulative-sum
matrix $C_H$. The level covariance is
$C_H\widehat V_\Delta C_H'$. For weak-IV inversion, first transform
the numerator moments and their full covariance, then invert; summing
separate confidence endpoints discards their dependence.

This route does not require estimating the cointegrating rank to define
the fixed-horizon functional. Efficient VECM fitting can use rank
restrictions as a challenger, with rank uncertainty retained. The state-space
connection is also substantive. In $T_{t+1}=T_t+\eta_{t+1}$,
$Y_{t+1}=BT_t+\epsilon_{t+1}$, the observed law identifies combinations
such as $B\Var(\eta)B'$, while invertible rotations of the latent trend
remain observationally equivalent. \citet[Theorems 1--2]{johansen2017}
shows that extracting the identified common-trend contribution and
estimating an individually normalized trend have different rate conditions.
Naming a filtered trend ``funding demand'' does not identify its intervention.

Lag augmentation offers a different persistence treatment. For
$Y_t=\phi Y_{t-1}+u_t$, regress $Y_{t+h}$ on $(Y_t,Y_{t-1})$.
The residualized first regressor approximates $u_t$ and the target
coefficient is $\phi^h$. Under mean independence of $u_t$ from the
innovations at \emph{all other dates}, the score formed with the
future projection error is serially uncorrelated even when that error
overlaps. \citet[Definition 1, Assumptions 1--3, Proposition 1]{lpaug2020}
proves marginal uniform inference for a specified VAR parameter space,
with eighth moments, summable innovation-product cumulants and scaled
design nonsingularity. Near unit roots its horizon condition is
$h_T/T\to0$. This theorem does not itself establish external-IV ratio
inference or simultaneous bands. We select the stationary-change IV route
for the integrated structural response and retain lag augmentation as
the justified alternative for its stated reduced-form projection target.

A regime break is a third problem. Pooling two different $F$ or impact
columns can average two economic mechanisms. A stationary transformation
does not remove a changing causal response. Segmented analyses require
enough dates in each segment, a declared treatment of estimated break
locations and, if used, a separate selection-aware uncertainty argument.
The product must identify that unresolved regime instead of interpreting
one pooled curve as a stable banking mechanism.
